% TOP_PROBLEM: {"id":30007523,"problem_number":"LOCAL-30007523","title":"Phillips-curve nonlinearities","category_id":21,"category":{"id":21,"name":"economics","display_name":"Economics","description":"Open economic theory statements, published research agendas, and proposed scoped specifications; question provenance and historical priority are qualified per record.","slug":"economics","order_index":21,"created_at":"2026-10-08T00:00:00Z"},"status":"research_question","external_url":null,"legacy_ids":[],"published":false,"statement":"Identify whether a causal inflation–slack relation has a stable threshold or slope change after accounting for supply shocks and measurement error.","economics":{"original_id":12,"field":"Inflation and monetary policy","kind":"identification","formalization_basis":"adapted_research_question","source_status":"research_agenda","priority_status":"not_established","priority_note":"Original priority is not established. No dated primary statement of this precise problem was verified.","earliest_verified_reference":{"authors":["European Central Bank"],"title":"Research agenda: Forecasting, macroeconomic and financial analysis","year":null,"url":"https://www.ecb.europa.eu/pub/economic-research/research_agenda/forecasting/html/index.en.html","doi":null,"locator":"“Understanding the inflation process,” selected areas of focus","evidence_summary":"The agenda explicitly includes Phillips-curve determinants and their evolution, wage–price dynamics, and the measurement, determinants, and role of inflation expectations."},"current_reference":{"authors":["European Central Bank"],"title":"Research agenda: Forecasting, macroeconomic and financial analysis","year":null,"url":"https://www.ecb.europa.eu/pub/economic-research/research_agenda/forecasting/html/index.en.html","doi":null,"locator":"“Understanding the inflation process,” selected areas of focus","evidence_summary":"The agenda explicitly includes Phillips-curve determinants and their evolution, wage–price dynamics, and the measurement, determinants, and role of inflation expectations."},"formalization_note":"A specific semiparametric identification problem is adapted from the ECB agenda. Neither this equation nor a universal Phillips-curve threshold is claimed to be published."},"statusQualification":"Published question or research agenda verified at the cited locator. The mathematical specification is adapted for this catalogue and includes additional assumptions; source publication does not establish first-ever priority or certify this exact model remains unsolved. A specific semiparametric identification problem is adapted from the ECB agenda. Neither this equation nor a universal Phillips-curve threshold is claimed to be published.","statusReviewedAt":"2026-10-08"}
% FIRST_POSTED: 2026-10-08
% ENTRY_KIND: statement_only
% !TeX program = lualatex
\documentclass[11pt]{article}
\usepackage{fontspec}
\usepackage[a4paper,margin=25mm]{geometry}
\usepackage{amsmath,amssymb,mathtools}
\usepackage{xurl}
\usepackage[hidelinks]{hyperref}
\newcommand{\sref}[2]{\hyperlink{source:#1:#2}{[S#2]}}
\setlength{\emergencystretch}{3em}
\begin{document}
% problemId:30007523
\section*{Phillips-curve nonlinearities}\label{30007523}
\subsection{Definitions and mathematical statement}
Let \(\pi_t\) be quarterly inflation, \(E_t^s\pi_{t+1}\) a specified survey measure of expectations, and \(x_t\) an explicitly constructed measure of real economic slack. Model the conditional inflation relation as
\[
\pi_t=\alpha+\beta E_t^s\pi_{t+1}
+\kappa_1 x_t+\kappa_2(x_t-\tau)_+
+\gamma'w_t+u_t,
\]
where \((r)_+=\max\{r,0\}\), \(w_t\) contains predetermined supply controls, \(\tau\) is a possible threshold, and \(u_t\) contains unobserved cost disturbances. Specify instruments \(Z_t\) satisfying \(E[Z_tu_t]=0\), relevance conditions, and the sample population before estimation.
The target is the identified set for \((\kappa_1,\kappa_2,\tau)\), allowing measurement error in slack and expectations and weak threshold identification when \(\kappa_2\) is near zero. Determine whether a slope change survives changes in the slack measure, supply controls, and monetary-policy regime. Compare a threshold model with smooth nonlinear alternatives on held-out observations; do not infer a kink solely from extreme inflation episodes that also experienced supply disruptions. State the invariance assumptions under which an estimated curve predicts inflation after a policy intervention. The open empirical agenda concerns the existence, location, and policy relevance of nonlinearities in a defined setting, rather than whether inflation and activity ever covary.

\par\textbf{Formulation and scope.} A specific semiparametric identification problem is adapted from the ECB agenda. Neither this equation nor a universal Phillips-curve threshold is claimed to be published.

\par\textbf{Open-question provenance.} An explicit question or research agenda was verified at the source locator. This does not by itself certify that every later version remains unresolved \sref{30007523}{1}.

\par\textbf{Historical priority.} Original priority is not established. No dated primary statement of this precise problem was verified.

\subsection{Short English statement}
Identify whether a causal inflation–slack relation has a stable threshold or slope change after accounting for supply shocks and measurement error.

\subsection{Sources}
\begin{itemize}\raggedright
\item[\textnormal{[S1]}]\hypertarget{source:30007523:1}{} European Central Bank. Year not verified. Research agenda: Forecasting, macroeconomic and financial analysis. “Understanding the inflation process,” selected areas of focus.
\url{https://www.ecb.europa.eu/pub/economic-research/research_agenda/forecasting/html/index.en.html}.
{\small\textit{Earliest verified question/agenda reference; historical priority qualified below. The agenda explicitly includes Phillips-curve determinants and their evolution, wage–price dynamics, and the measurement, determinants, and role of inflation expectations.}}
\end{itemize}
\end{document}
